\chapter{}

Agostino, figliolo di un cavallaio che aveva due poderi a confine con Poggio a’ Meli, non voleva che Pietro parlasse troppo a Ghìsola; per quell’amor proprio che nell’adolescenza somiglia alla gelosia. E capì che doveva odiare il rispetto ingenuo di Pietro; e compatirlo come una debolezza.

Ghìsola, infatti, dava al suo padroncino un senso di disagio e d’impaccio; ma egli voleva essere forte e cercava di convincersi che preferiva l’amicizia di Agostino; e con lui doventava remissivo ed obbediente; procurando d’indovinare le cose che pensava e non diceva a posta. Talvolta gli raccattava una pietra com’egli comandava soltanto guardandola; per tirarla a pena visto un uccello sopra un ramo accanto alla strada. E come il vento gonfiava la camicia d’Agostino, tutta sbottonata! Perchè non aveva i polsi eguali a lui, le ciglia, gli orecchi, la camicia? E perchè quando si provava a fare come lui, con la stessa aria di noncuranza, si trovava perso d’animo, senza fiato, con la paura di provocare la sua collera che lo faceva tremare? Perchè non poteva sostenere il suo sguardo crucciato, impenetrabile e lucido, quando si provava a non rispondere alle sue domande e quando non aveva indovinato? Quello sguardo lo impauriva così come quando, senza essercene avvisti prima, ci si trova proprio ai piedi una fonte piena d’acqua.

Agostino aveva il naso piccolo e corto, di bambino, tutto lentigginoso; ma il suo collo era come quello di una bella donna; le mani fatte bene. I suoi colloqui con Ghìsola, che consistevano in parole senza senso, convenzionali, che capivano loro due soltanto, suscitavano in Pietro sentimenti inaspettati; ai quali da solo non avrebbe mai sognato. E il diletto d’ascoltarli era tanto! Anche gli pareva d’imparare chi sa che.

Ghìsola aveva un sorriso piacevole dicendo certe cose, che a lei sola potevano venire in mente; e Pietro si struggeva dalla voglia d’impararle come i suoi stornelli. Ma non riusciva nè meno a cantare; e ne aveva vergogna. Talvolta non volendo che ridesse, le faceva qualche dispetto a posta.

Sotto il largo cappello di paglia, che le calava sempre sopra un orecchio, guarnito con un nastro di raso liso e con due rosette buttate via da Anna, il volto di Ghìsola era tranquillamente insignificante e sciatto.

Sembrava, con la sottana rimendata male, troppo semplice e quasi stupida.

Vi sono esseri che non chiedono nulla a nessuno e rinunziano a tutto; e, non essendo rispettati come gli altri, pare che di loro se ne possa fare quel che si vuole. Perciò quel che riguarda gli altri lo trovano antipatico. Se qualcuno li ama, non vogliono cambiarsi; chiedendo che cosa questo bene esiga. E allora lo evitano.

Quando Masa batteva le nocche su la fronte di Ghìsola dicendole: «che ci hai qui?» ella rispondeva quasi con esasperazione:

— Che ne sapete voi? Che ve ne importa?

Talvolta credeva, con piacere e con stizza, che il suo viso offendesse. Quando gli altri parlavano si metteva silenziosa; credendoli diffidenti. Non la interessava niente; obbediva a Masa e ai padroni, perchè da sè stessa non avrebbe pensato nè meno alla calza; e sentiva malvolentieri che tutto ciò che esiste non era soltanto in lei.

Talvolta pareva piuttosto che parlasse con lo scalone di casa; quando, secondo il suo solito, ci stava seduta.

Non si sarebbe arrischiata ad avere qualche idea, perchè ne aveva troppe che non le si addicevano; come non si arrischiava, quando era andata alla trattoria, a chiedere le ghiottonerie che vedeva; e invece le avvampavano il viso, e la stordivano quanto le stanze calde a cui non era abituata.

Ma c’era in lei il presentimento e il senso di una vita, che le montava la testa come la ricchezza e il lusso degli altri.
